\documentclass[10pt]{article}
\usepackage[a4paper,margin=20mm]{geometry}
\usepackage[T1]{fontenc}
\usepackage{lmodern}
\usepackage{microtype}
\usepackage{amsmath,amssymb,mathtools,bm}
\usepackage{booktabs,tabularx,array}
\usepackage[dvipsnames]{xcolor}
\usepackage[hidelinks]{hyperref}
\usepackage{enumitem}
\setlist{nosep,leftmargin=1.5em}
\setlength{\parindent}{0pt}
\setlength{\parskip}{5pt}
\definecolor{navy}{HTML}{183B56}
\definecolor{teal}{HTML}{127475}
\definecolor{soft}{HTML}{F2F6F7}
\newcommand{\R}{\mathbb{R}}
\newcommand{\C}{\mathbb{C}}
\newcommand{\wrap}{\operatorname{wrap}}
\newcommand{\LN}{\operatorname{LN}}
\newcommand{\MLP}{\operatorname{MLP}}
\newcommand{\maskP}{\mathcal M_P}
\newcommand{\maskQ}{\mathcal M_Q}
\newcommand{\code}[1]{\texttt{#1}}
\title{\vspace{-1.4em}\color{navy}\textbf{Released GridFM GraphKit as an AC Power-Flow Surrogate}\\
\large Exact formulation of the PPC branch-row adaptation}
\author{Implementation note for \code{train\_valid\_test\_gridfm.py}}
\date{August 2026}

\begin{document}
\maketitle
\vspace{-1.5em}

\begin{center}
\fcolorbox{teal}{soft}{\parbox{0.93\linewidth}{
\textbf{Scope.} This document formulates the released
\code{GNS\_heterogeneous} from the released \code{gridfm\_graphkit} package
used with \code{--task pf --gridfm\_impl graphkit}. It is not the earlier local
\code{GridFMHeteroSurrogate} mirror. The optional post-model known-operator
projection is excluded; all equations describe the raw GraphKit prediction.
}}
\end{center}

\section{AC Power-Flow Problem}

Let bus voltage be
\begin{equation}
  \underline V_i = v_i e^{\mathrm j\theta_i},\qquad
  \underline{\bm V}=[\underline V_1,\ldots,\underline V_N]^\top.
\end{equation}
For the PPC-derived admittance matrix \(\bm Y\), calculated nodal power is
\begin{equation}
  \underline{\bm S}^{\rm calc}(\bm v,\bm\theta)
  =\underline{\bm V}\odot\overline{\bm Y\underline{\bm V}}.
\end{equation}
With specified injection \(\underline{\bm S}^{\rm set}=\bm P^{\rm set}+\mathrm j\bm Q^{\rm set}\),
\begin{equation}
  \Delta\bm P=\bm P^{\rm set}-\Re(\underline{\bm S}^{\rm calc}),\qquad
  \Delta\bm Q=\bm Q^{\rm set}-\Im(\underline{\bm S}^{\rm calc}).
\end{equation}
The PF equation masks are
\begin{equation}
 \maskP=\{i:i\notin\mathcal V_{\rm REF}\},\qquad
 \maskQ=\{i:i\in\mathcal V_{\rm PQ}\}.
\end{equation}
Thus active balance is imposed at PV and PQ buses, while reactive balance is
imposed only at PQ buses.

\section{Adapted Heterogeneous Graph}

The adapted graph is
\begin{equation}
 \mathcal G=(\mathcal V_b\cup\widetilde{\mathcal V}_g,
 \mathcal E_{bb}\cup\mathcal E_{gb}\cup\mathcal E_{bg}).
\end{equation}
Here \(\mathcal V_b\) contains one node per PPC bus. In native GridFM,
\(\mathcal V_g\) contains physical generator rows. The PPC PF parquet has no
generator-resolved table, so this adapter constructs one \emph{proxy} generator
\(\tilde g_i\) for each PV or REF bus \(i\):
\begin{equation}
 b(\tilde g_i)=i,\qquad
 |\widetilde{\mathcal G}(i)|=
 \begin{cases}1,&i\in\mathcal V_{\rm PV}\cup\mathcal V_{\rm REF},\\0,&i\in\mathcal V_{\rm PQ}.\end{cases}
\end{equation}
The two association edges \(\tilde g_i\to i\) and \(i\to\tilde g_i\) do not
represent electrical branches; they allow two-way latent exchange.

\subsection{Bus and proxy-generator inputs}

Define \(T(x)=\operatorname{sign}(x)\log(1+|x|)\). The bus row is
\begin{align}
 \bm x_i^b=[&T(-P_i^{\rm set}),T(-Q_i^{\rm set}),0,
 v_i^{0},\theta_i^{0},\mathbb 1_{\rm PQ},\mathbb 1_{\rm PV},\mathbb 1_{\rm REF},\nonumber\\
 &0.5,1.5,-10,10,T(G_i^{\rm sh}),T(B_i^{\rm sh}),\log_{10}(V_{n,i}^{\rm kV}+10^{-9})]\in\R^{15}.
\end{align}
The standard PF adapter uses
\begin{equation}
 \bm x_{\tilde g_i}^g=
 [T(\max(P_i^{\rm set},0)),-10,10,0,0,0]\in\R^6.
\end{equation}
The limits and costs are placeholders, not reconstructed physical-generator
data. The voltage start \((v^0,\theta^0)\) is retained in the bus row.

\subsection{Transformer-aware directed branch features}

For branch \(f\to t\), let
\begin{equation}
 a=\tau e^{\mathrm j\phi},
\end{equation}
and use the already per-unit branch quantities stored by the PPC loader:
\begin{align}
 Y_{ff}&=\frac{y_f+y_f^{\rm sh}/2}{|a|^2}, &
 Y_{ft}&=-\frac{y_f}{\overline a},\\
 Y_{tt}&=y_t+y_t^{\rm sh}/2, &
 Y_{tf}&=-\frac{y_t}{a}.
\end{align}
Each physical branch becomes two directed message edges. Their features are
\begin{align}
 \bm e_{ft}&=[0,0,\Re Y_{ff},\Im Y_{ff},\Re Y_{ft},\Im Y_{ft},\tau,0,0,S^{\max}],\\
 \bm e_{tf}&=[0,0,\Re Y_{tt},\Im Y_{tt},\Re Y_{tf},\Im Y_{tf},\tau,0,0,S^{\max}].
\end{align}
The first two flow channels and the angle-limit channels are zero in the current
PF adapter; the admittance split and tap carry the meaningful electrical data.
The adapter's reconstruction validator confirms
\begin{equation}
 \bm Y_{\rm adapter}=\bm Y_{\rm parquet}
\end{equation}
to floating-point tolerance, including bus shunts.

\section{Relation-Specific Transformer Layers}

Initial embeddings are produced independently by node and edge type:
\begin{equation}
 \bm h_i^{b,(0)}=E_b(\bm x_i^b),\quad
 \bm h_g^{g,(0)}=E_g(\bm x_g^g),\quad
 \bar{\bm e}_{ji}=E_e(\bm e_{ji}),
\end{equation}
where each \(E\) is Linear--LeakyReLU--Linear--LayerNorm. The released
configuration uses hidden width \(d=48\), \(H=8\) heads, and \(L=12\) layers.

At layer \(\ell\), each relation has its own TransformerConv operator:
\begin{equation}
 \mathcal T_{bb}^{(\ell)},\qquad
 \mathcal T_{gb}^{(\ell)},\qquad
 \mathcal T_{bg}^{(\ell)}.
\end{equation}
For a bus edge \(j\to i\) and head \(h\), a faithful operator-level form is
\begin{align}
 \bm q_i^{h}&=W_{Q,bb}^{h}\bm h_i^b,\\
 \bm k_{ji}^{h}&=W_{K,bb}^{h}\bm h_j^b+W_{K,e}^{h}\bar{\bm e}_{ji},\\
 \bm z_{ji}^{h}&=W_{V,bb}^{h}\bm h_j^b+W_{V,e}^{h}\bar{\bm e}_{ji},\\
 \alpha_{ji}^{h}&=\operatorname{softmax}_{j\in\mathcal N_{bb}(i)}
 \left(\frac{(\bm q_i^h)^\top\bm k_{ji}^h}{\sqrt{d_h}}\right),\\
 \bm m_{bb,i}^{(\ell)}&=W_{O,bb}\mathop{\Vert}_{h=1}^{H}
 \sum_{j\in\mathcal N_{bb}(i)}\alpha_{ji}^{h}\bm z_{ji}^{h}.
\end{align}
Association relations are analogous but have no branch attribute:
\begin{align}
 \bm m_{gb,i}^{(\ell)}&=\mathcal T_{gb}^{(\ell)}
 (\bm h_i^b,\{\bm h_g^g:b(g)=i\}),\\
 \bm m_{bg,g}^{(\ell)}&=\mathcal T_{bg}^{(\ell)}
 (\bm h_g^g,\bm h_{b(g)}^b).
\end{align}
HeteroConv sums relation outputs at each target type, then applies
LayerNorm, LeakyReLU, and a shape-compatible residual connection:
\begin{align}
 \bm h_i^{b,(\ell+1)}&\leftarrow \bm h_i^{b,(\ell)}+
 \sigma\!\left(\LN(\bm m_{bb,i}^{(\ell)}+\bm m_{gb,i}^{(\ell)})\right),\\
 \bm h_g^{g,(\ell+1)}&\leftarrow \bm h_g^{g,(\ell)}+
 \sigma\!\left(\LN(\bm m_{bg,g}^{(\ell)})\right).
\end{align}
The first layer changes width from \(48\) to \(48H\), so its outer residual is
omitted; subsequent layers retain it. TransformerConv also uses its internal
\(\beta\)-gated root/neighbor combination.

\section{Layerwise Physics Feedback}

Unlike the local mirror, the released model invokes the same bus and generator
decoders after \emph{every} layer:
\begin{equation}
 [\tilde v_i^{(\ell)},\tilde\theta_i^{(\ell)}]
 =D_b(\bm h_i^{b,(\ell)}),\qquad
 \tilde P_g^{(\ell)}=D_g(\bm h_g^{g,(\ell)}).
\end{equation}
PF masks pin known state components:
\begin{align}
 v_i^{(\ell)}&=
 \begin{cases}\tilde v_i^{(\ell)},&i\in\mathcal V_{\rm PQ},\\v_i^0,&i\in\mathcal V_{\rm PV}\cup\mathcal V_{\rm REF},\end{cases}\\
 \theta_i^{(\ell)}&=
 \begin{cases}\theta_i^0,&i\in\mathcal V_{\rm REF},\\\tilde\theta_i^{(\ell)},&\text{otherwise}.
 \end{cases}
\end{align}

For a directed edge \(i\to j\), the branch-current and branch-power operators are
\begin{align}
 \underline I_{ij}^{(\ell)}&=Y_{ii}^{(ij)}\underline V_i^{(\ell)}+
 Y_{ij}^{(ij)}\underline V_j^{(\ell)},\\
 \underline S_{ij}^{(\ell)}&=\underline V_i^{(\ell)}
 \overline{\underline I_{ij}^{(\ell)}}.
\end{align}
Outgoing directed powers are aggregated at each bus:
\begin{equation}
 P_i^{\rm net,(\ell)}=\sum_{j:(i,j)\in\mathcal E_{bb}}P_{ij}^{(\ell)},\qquad
 Q_i^{\rm net,(\ell)}=\sum_{j:(i,j)\in\mathcal E_{bb}}Q_{ij}^{(\ell)}.
\end{equation}
The PF physics decoder analytically recovers quantities that are not specified:
\begin{align}
 Q_{g,i}^{(\ell)}&=
 \begin{cases}
 Q_i^{\rm net,(\ell)}+Q_{d,i}-q_i^{\rm sh,(\ell)},
 &i\in\mathcal V_{\rm PV}\cup\mathcal V_{\rm REF},\\0,&i\in\mathcal V_{\rm PQ},
 \end{cases}\\
 P_{g,i}^{(\ell)}&=
 \begin{cases}
 \sum_{g:b(g)=i}\tilde P_g^{(\ell)},&i\in\mathcal V_{\rm PV},\\
 P_i^{\rm net,(\ell)}+P_{d,i}-p_i^{\rm sh,(\ell)},&i\in\mathcal V_{\rm REF},\\
 0,&i\in\mathcal V_{\rm PQ}.
 \end{cases}
\end{align}
With \(p_i^{\rm sh}=-G_i^{\rm sh}v_i^2\) and
\(q_i^{\rm sh}=B_i^{\rm sh}v_i^2\), the layer residual is
\begin{align}
 r_{P,i}^{(\ell)}&=P_{g,i}^{(\ell)}-P_{d,i}+p_i^{\rm sh,(\ell)}-P_i^{\rm net,(\ell)},\\
 r_{Q,i}^{(\ell)}&=Q_{g,i}^{(\ell)}-Q_{d,i}+q_i^{\rm sh,(\ell)}-Q_i^{\rm net,(\ell)}.
\end{align}
It is injected into the next bus representation:
\begin{equation}
 \bm h_i^{b,(\ell)}\leftarrow \bm h_i^{b,(\ell)}+
 \MLP_{\rm phys}([r_{P,i}^{(\ell)},r_{Q,i}^{(\ell)}]).
\end{equation}
Only the last layer's decoded state is returned. This is depthwise learned
physics feedback, not a Newton step and not a \(K\)-step voltage unrolling.

\section{Output and Training Objective}

The raw AC-PF prediction is
\begin{equation}
 \widehat{\bm V}_{\rm state}=
 \{(\hat v_i,\hat\theta_i)=(v_i^{(L)},\wrap(\theta_i^{(L)}))\}_{i=1}^{N}.
\end{equation}
Against Newton labels \((v_i^\star,\theta_i^\star)\), the supervised term is
\begin{equation}
 \mathcal L_{\rm state}=\frac1N\sum_i
 \left[(\hat v_i-v_i^\star)^2+wrap(\hat\theta_i-\theta_i^\star)^2\right].
\end{equation}
The external PPC physics term recomputes power from the exact dataset \(\bm Y\),
not from transformed neural features:
\begin{equation}
 \mathcal L_{\rm phys}=\frac1{|\maskP|+|\maskQ|}
 \left[\sum_{i\in\maskP}\rho(\Delta P_i)+
 \sum_{i\in\maskQ}\rho(\Delta Q_i)\right],
\end{equation}
where the default robust penalty is
\begin{equation}
 \rho(r)=\log\cosh(\operatorname{clip}(r,-30,30)).
\end{equation}
The default total objective is
\begin{equation}
 \boxed{\mathcal L=\mathcal L_{\rm state}+10^{-2}\mathcal L_{\rm phys}.}
\end{equation}
Training physics is evaluated through the script's complex64 path; reported
validation/test residual metrics reconstruct \(\underline{\bm S}^{\rm calc}\)
with complex128 when the dataset is loaded with
\code{--dataset\_complex\_dtype complex128}.

\section{Interpretive Summary}

\begin{tabularx}{\linewidth}{>{\bfseries}p{0.25\linewidth}X}
\toprule
Property & Released GraphKit PF adaptation \\
\midrule
Graph & Bus nodes plus one proxy generator per PV/REF bus; three directed relation types.\\
Electrical edge encoding & Directed transformer-aware \(\pi\)-model split \((Y_{ff},Y_{ft})\) and \((Y_{tt},Y_{tf})\).\\
Message passing & 12 relation-specific multi-head TransformerConv layers, width 48 and 8 heads.\\
Voltage decoding & Shared decoder at every layer; known PF quantities are pinned.\\
Physics inside forward & Branch flows, nodal injections, analytic PF decoder, and residual feedback at every layer.\\
Physics outside forward & Exact-Ybus robust residual loss and complex128 validation/test metrics.\\
Pretraining status & No released GridFM checkpoint; this implementation is trained from scratch.\\
\bottomrule
\end{tabularx}

\section*{Code provenance}
{\small\sloppy
The equations follow \path{gridfm_graphkit/models/gnn_heterogeneous_gns.py},
\path{gridfm_graphkit/models/utils.py}, \path{gridfm_graphkit_adapter.py},
and the AC-PF path in \path{train_valid_test_gridfm.py}, as present in the
local workspace in August 2026.
\par}

\end{document}
